\documentclass[11pt,letterpaper]{article}
\usepackage[margin=1in]{geometry}
\usepackage{graphicx} % Required for inserting images
\usepackage{tikz}
\usepackage{amsthm,amsmath,amssymb}
\usepackage{dsfont}
\usepackage{algorithm}
\usepackage{algpseudocode}
\usepackage{xcolor}
\usepackage{soul}
\usepackage{enumitem}

\usepackage[hidelinks]{hyperref}
\usepackage{cleveref}
\usepackage{etoolbox}

% Theorems
\newtheorem{theorem}{Theorem}[section]
\newtheorem{corollary}[theorem]{Corollary}
\newtheorem{lemma}[theorem]{Lemma}
\newtheorem{claim}[theorem]{Claim}
\theoremstyle{definition}
\newtheorem{definition}[theorem]{Definition}
\newtheorem{remark}[theorem]{Remark}
\newtheorem{observation}[theorem]{Observation}

\crefname{theorem}{Theorem}{Theorems}
\Crefname{theorem}{Theorem}{Theorems}
\crefname{corollary}{Corollary}{Corollaries}
\Crefname{corollary}{Corollary}{Corollaries}
\crefname{lemma}{Lemma}{Lemmas}
\Crefname{lemma}{Lemma}{Lemmas}
\crefname{claim}{Claim}{Claims}
\Crefname{claim}{Claim}{Claims}
\crefname{definition}{Definition}{Definitions}
\Crefname{definition}{Definition}{Definitions}
\crefname{remark}{Remark}{Remarks}
\Crefname{remark}{Remark}{Remarks}
\crefname{observation}{Observation}{Observations}
\Crefname{observation}{Observation}{Observations}

% The environments above share the `theorem` counter, and cleveref names a label
% after its counter -- so without these aliases every \cref to a lemma, claim,
% etc. comes out as "Theorem". \crefalias is local to the environment's group,
% so each one is typed correctly while the shared numbering is left alone.
\AtBeginEnvironment{corollary}{\crefalias{theorem}{corollary}}
\AtBeginEnvironment{lemma}{\crefalias{theorem}{lemma}}
\AtBeginEnvironment{claim}{\crefalias{theorem}{claim}}
\AtBeginEnvironment{definition}{\crefalias{theorem}{definition}}
\AtBeginEnvironment{remark}{\crefalias{theorem}{remark}}
\AtBeginEnvironment{observation}{\crefalias{theorem}{observation}}

\newcommand{\reals}{\mathbb{R}}
\newcommand{\Span}{\mathrm{span}}
\newcommand{\coreq}{\mathrm{coreq}}
\newcommand{\cone}{\mathrm{cone}}

\newcommand{\argmax}{\operatorname*{arg\,max}}
\newcommand{\argmin}{\operatorname*{arg\,min}}

\newcommand{\calA}{\mathcal{A}}
\newcommand{\calB}{\mathcal{B}}
\newcommand{\calC}{\mathcal{C}}
\newcommand{\calD}{\mathcal{D}}
\newcommand{\calE}{\mathcal{E}}
\newcommand{\calF}{\mathcal{F}}
\newcommand{\calG}{\mathcal{G}}
\newcommand{\calH}{\mathcal{H}}
\newcommand{\calI}{\mathcal{I}}
\newcommand{\calJ}{\mathcal{J}}
\newcommand{\calK}{\mathcal{K}}
\newcommand{\calL}{\mathcal{L}}
\newcommand{\calM}{\mathcal{M}}
\newcommand{\calN}{\mathcal{N}}
\newcommand{\calO}{\mathcal{O}}
\newcommand{\calP}{\mathcal{P}}
\newcommand{\calQ}{\mathcal{Q}}
\newcommand{\calR}{\mathcal{R}}
\newcommand{\calS}{\mathcal{S}}
\newcommand{\calT}{\mathcal{T}}
\newcommand{\calU}{\mathcal{U}}
\newcommand{\calV}{\mathcal{V}}
\newcommand{\calW}{\mathcal{W}}
\newcommand{\calX}{\mathcal{X}}
\newcommand{\calY}{\mathcal{Y}}
\newcommand{\calZ}{\mathcal{Z}}

\newcommand{\bfa}{\mathbf{a}}
\newcommand{\bfb}{\mathbf{b}}
\newcommand{\bfc}{\mathbf{c}}
\newcommand{\bfd}{\mathbf{d}}
\newcommand{\bfe}{\mathbf{e}}
\newcommand{\bff}{\mathbf{f}}
\newcommand{\bfg}{\mathbf{g}}
\newcommand{\bfh}{\mathbf{h}}
\newcommand{\bfi}{\mathbf{i}}
\newcommand{\bfj}{\mathbf{j}}
\newcommand{\bfk}{\mathbf{k}}
\newcommand{\bfl}{\mathbf{l}}
\newcommand{\bfm}{\mathbf{m}}
\newcommand{\bfn}{\mathbf{n}}
\newcommand{\bfo}{\mathbf{o}}
\newcommand{\bfp}{\mathbf{p}}
\newcommand{\bfq}{\mathbf{q}}
\newcommand{\bfr}{\mathbf{r}}
\newcommand{\bfs}{\mathbf{s}}
\newcommand{\bft}{\mathbf{t}}
\newcommand{\bfu}{\mathbf{u}}
\newcommand{\bfv}{\mathbf{v}}
\newcommand{\bfw}{\mathbf{w}}
\newcommand{\bfx}{\mathbf{x}}
\newcommand{\bfy}{\mathbf{y}}
\newcommand{\bfz}{\mathbf{z}}

\newcommand{\cost}{\mathrm{COST}}
\newcommand{\ALG}{\mathrm{ALG}}
\newcommand{\OPT}{\mathrm{OPT}}
\newcommand{\Flow}{\textsc{Flow}}
\newcommand{\ouralg}{\textsc{LineTracker}}

\newcommand{\EE}{\mathbb{E}}
\newcommand{\Var}{\mathrm{Var}}
\newcommand{\Conv}{\mathrm{Conv}}
\newcommand{\One}{\mathds{1}}

\newcommand{\interval}[1]{\lfloor #1 \rceil}
\newcommand{\dis}[1]{#1}
\newcommand{\off}{\mathrm{off}}
\newcommand{\on}{\mathrm{on}}
\newcommand{\org}{\mathrm{org}}
\newcommand{\rel}{\mathrm{rel}}
\newcommand{\new}{\mathrm{new}}

\newcommand{\solution}{\medskip\noindent\textbf{Solution.}\par\medskip}

\title{47-851: Modern Convex Optimization \\ Homework 1}
\author{Poon Tze-Yang}
\date{November 2026}


\begin{document}

\maketitle

\noindent Lecturer: Karan Singh

\begin{enumerate}[series=q]
\item \textbf{Inner Products and Dual Norms (10 points)}
\begin{enumerate}[label=\textbf{Part \Alph*.}, leftmargin=*]
\item Prove that a norm on a real vector space $V$ is induced by an inner product if and only if it satisfies the parallelogram identity
\[
\|x + y\|^2 + \|x - y\|^2 = 2\|x\|^2 + 2\|y\|^2 \quad \forall x, y \in V.
\]
\item For a norm $\|\cdot\|$ on $\reals^n$, define its dual norm as
\[
\|x\|_* = \max_{\|y\| \leq 1} y^\top x.
\]
\begin{enumerate}[label=\arabic*.]
\item Calculate the dual norms of $\|\cdot\|_2$ and $\|\cdot\|_\infty$, where $\|x\|_2 = \sqrt{x^\top x}$ and $\|x\|_\infty = \max_i |x_i|$.
\item Show that
\[
\partial \|\cdot\|(x) = \{y : \|y\|_* \leq 1, \ y^\top x = \|x\|\}.
\]
\end{enumerate}
\end{enumerate}
\end{enumerate}

\solution
(A)

$(\Rightarrow):$
\begin{align*}
 \left\| x+y\right\|^2 + \left\|x-y\right\| ^2 & = \left\langle x+y,x+y\right\rangle + \left\langle x-y,x-y\right\rangle \\
 & = 2\left\langle x,x\right\rangle + 2\left\langle y,y\right\rangle + (\left\langle x,y\right\rangle + \left\langle y,x\right\rangle) - (\left\langle x,y\right\rangle + \left\langle y,x\right\rangle) \\
 & = 2\left\|x\right\|^2 + 2\left\|y\right\|^2.
\end{align*}

$(\Leftarrow):$
HARD, READ MORE

Given a norm $\|\cdot\|$ that satisfies the parallelogram identity, we define $\langle\cdot,\cdot\rangle : V\times V \mapsto \reals$ by 
\begin{align*}
    \langle x,y\rangle = \frac{1}{2}\left(\|x+y\|^2 - \|x\|^2 - \|y\|^2\right).
\end{align*}
\begin{lemma}\label{lemma:sandwichNorm}

\end{lemma}
\begin{proof}
    From the definition of a norm, we have that 
    \begin{align*}
        \|x+y\|\leq \|x\|+\|y\| \ \Rightarrow \|x+y\|^2 \leq \|x\|^2 + \|y\|^2 +2\|x\|\|y\|
    \end{align*}
    \begin{align*}
        \|x\|\|y\| \leq \frac{1}{2}(\|x+y\|^2 - \|x\|^2 - \|y\|^2)
    \end{align*}
\end{proof}

To show that this defines an inner product, we verify the definition of an inner product:
\begin{enumerate} 
    \item \ul{$\langle x,x\rangle\geq 0\ \forall x$ and $\langle x,x\rangle = 0 $ iff $x=\mathbf{0}$}: Plugging $\|x+x\|^2 = 4\|x\|^2$ into the defintion of the inner product gives $\langle x,x\rangle = \|x\|^2$ and the result follows from the definition of a norm. 
    \item \ul{$\langle x,y\rangle = \langle y,x\rangle\ \forall x,y\in V$:} Follows since addition on $x+y$ is commutative.
    \item \ul{$\langle x, y+z\rangle = \langle x,y\rangle + \langle x,z\rangle \ \forall x,y,z\in V$:}
    \begin{align*}
        \langle x,y+z\rangle &= \frac{1}{4}(\|x+y+z\|^2 - \|x-y-z\|^2)\\
        &= \frac{1}{4}(-\|x+z-y\|^2 + 2\|x+z\|^2+2\|y\|^2 + \|x-y+z\|^2 - 2\|x-y\|^2-2\|z\|^2)\\
        &=\frac{1}{4}(\|x+z\|^2-\|x-y\|^2) + \frac{1}{4}(\|x+z\|^2-2\|z\|^2 -\|x-y\|^2+2\|y\|^2) \\
        &=\frac{1}{4}(\|x+z\|^2-\|x-y\|^2) + \frac{1}{4}(\|x+y\|^2-\|x-z\|^2 +2\|x\|^2-2\|x\|^2) \\
        &=\frac{1}{4}(\|x+y\|^2-\|x-y\|^2) + \frac{1}{4}(\|x+z\|^2-\|x-z\|^2) \\
        &=\langle x,y\rangle + \langle x,z\rangle.
    \end{align*}
    \item \ul{$\langle \alpha x,y\rangle = \alpha\langle x,y\rangle\ \forall \alpha\in \reals,\ \forall x,y\in V$:}
    From the additivity property above, homogeneity holds for $\alpha \in \mathbb{Z}$. For any rational $\alpha = p/q$ for integers $p,q$, and letting $qz := x$ we have that
    \[
        \langle \frac{p}{q} x,y\rangle = \langle p z, y\rangle = p\langle z,y\rangle  = \frac{p}{q}\langle qz,y\rangle = \frac{p}{q}\langle x,y\rangle,
    \]
    so homogeneity holds over $\mathbb{Q}$.
    \begin{lemma} [Cauchy-Schwarz holds]
        $\langle x,y\rangle \leq \|x\|\|y\|$.
    \end{lemma}
    \begin{proof}
        \[\langle x,y\rangle = \frac{1}{2}\left(\|x+y\|^2 - \|x\|^2 - \|y\|^2\right) \leq \frac{1}{2}\left(- \|x\|^2 - \|y\|^2\right) \leq \frac{1}{2}(2\|x\|\|y\|) = \|x\|\|y\|\]
    \end{proof}
    Finally, for $\lambda\in \reals$ and $r\in \mathbb{Q}$, since additivity holds for rationals, we have that 
    \begin{align}
        \langle \lambda x,y\rangle - \lambda\langle x,y\rangle &= \langle \lambda x,y\rangle - (\lambda-r)\langle x,y\rangle + r\langle x,y\rangle
    \end{align}
    Taking the absolute value on both sides and using Cauchy-Schwarz:
    \begin{align*}
        |\langle \lambda x,y\rangle - \lambda\langle x,y\rangle | &\leq |\langle \lambda x,y\rangle| + r|\langle x,y\rangle| + |\lambda-r||\langle x,y\rangle|\\
        & \leq 2|\lambda-r| \|x\| \|y\|.
    \end{align*}
    Since we can choose $r$ arbitrarily close to $\lambda$, we have that $|\langle \lambda x,y\rangle - \lambda\langle x,y\rangle |=0$, so homogeneity holds over reals.
\end{enumerate}

(B):

1.

$\|\cdot\|_2$:
By Cauchy-Schwarz, we can upper bound the value of the dual norm $y^\top x\leq \|x\|_2\|y\|_2\leq \|x\|_2$. 
Select $\hat{y}=\frac{x}{\|x\|_2}$, so $\|\hat{y}\|_2 = \sqrt{\frac{x^\top x}{\|x\|^2}} = 1$ so $\hat{y}$ is feasible. 
Furthermore, $\hat{y}^\top x = \|x\|_2$, which attains the upper bound and is therefore the value of the dual norm.
Hence, the dual norm of $\|\cdot\|_2$ is itself.

$\|\cdot\|_\infty$:
By definition of $\|\cdot\|_\infty$, we have:
\begin{align*}
    \max_{\|y\|_\infty\leq 1} y^\top x  = \sum_{i=1}^n\max_{|y_i|\leq 1} y_i x_i = \sum_{i=1}^n x_i.
\end{align*}
Hence, the dual norm of $\|\cdot\|_\infty$ is $\|\cdot\|_1$.

2.

By definition of subgradient:
\begin{align*}
    \partial\|\cdot \|(x) &= \{y:\|x'\|\geq\|x\| + (x'-x)^\top y\ \forall x'\in \reals^n\} \\
    &= \{y: y^\top x - \|x\|\geq y^\top x' - \| x'\|\ \forall x'\in \reals^n\}.
\end{align*}
For any $y\in \partial\|\cdot \|(x)$, note that $y^\top x - \|x\|$ is an upper bound on $y^\top x' - \| x'\|$. 
Suppose for contradiction there exists $x''\in \reals^n$ such that $y^\top x'' - \| x''\| = c>0$. 
Then $kx''$ gives $y^\top (kx'') - \| kx'\| = kc$ which exceeds $y^\top x - \|x\|$ as $k\to \infty$, which is a contradiction.
Thus, we have $y^\top x' - \| x'\|\leq 0 \ \forall x'\in \reals^n$.
In particular, $y^\top x - \| x\| \leq 0$. 
Furthermore, plugging in $x' = \mathbf{0}$ gives $y^\top x - \| x\| \geq 0$ and thus $y^\top x - \| x\| = 0$.
For every $\|x'\|\leq 1$, we have that $y^\top x'\leq \| x'\|\leq 1$, so $\left\|y\right\| = \max_{\left\|x'\right\|\leq 1}y^\top x'\leq 1$.
This gives the forward inclusion.

For the reverse inclusion, take any $y$ such that $\|y\|_* \leq 1, \ y^\top x = \|x\|$.
For all $x'\in \reals^n$, we show that $y^\top x' \leq \left\|x'\right\|$.
If $x'=\mathbf{0}$, this is true.
Otherwise, letting $\hat{x} = \frac{x'}{\|x'\|}$, since $\|\hat{x}\|\leq 1$, the condition is equivalent to $y^\top \hat{x}\leq 1$, which is true because $\|y\|_* =\max_{\|\hat{x}\|\leq 1}y^\top\hat{x}\leq 1$.
Since $y^\top x=\left\|x\right\|$ attains this bound, we obtain $y^\top x - \|x\|\geq y^\top x' - \| x'\|\ \forall x'\in \reals^n$, which gives the reverse inclusion.


\begin{enumerate}[resume=q]
\item \textbf{Local Descent and Checking PSD (10 points)}
\begin{enumerate}[label=\textbf{Part \Alph*.}, leftmargin=*]
\item Find a function $f : \reals^n \to \reals$ and a point $x \in \reals^n$ such that there is no (non-zero) descent direction at $x$, and yet $x$ is not a local minimizer.
\item Given $x \in \reals^n$, form a $(n + 1) \times (n + 1)$-dimensional matrix
\[
X = \begin{bmatrix}
1 & x_1 & x_2 & \cdots & x_n \\
x_1 & x_1 & 0 & \cdots & 0 \\
x_2 & 0 & x_2 & \cdots & 0 \\
\vdots & \vdots & \vdots & \ddots & \vdots \\
x_n & 0 & 0 & \cdots & x_n
\end{bmatrix}.
\]
Show that $X \succeq 0$ if and only if $x_i \geq 0, \forall i \in [n]$ and $\sum_{i=1}^n x_i \leq 1$.
\end{enumerate}
\end{enumerate}

\solution

A.

Select the function $f(x) = \sin\left(\frac{1}{\|x\|_2}\right)$, defined to be zero at $\mathbf{0}$.
At $x=\mathbf{0}$, there is no descent direction $d\neq \mathbf{0}$: for any $\varepsilon>0$, select a sufficiently integer large $k$ such that $\alpha =\frac{1}{(2k\pi + \frac{\pi}{2})\left\|d\right\|_2} < \varepsilon$.
This gives 
\[
f(\alpha d) = \sin\left(2k\pi + \frac{\pi}{2}\right) = 1>f(\mathbf{0}).
\]
Thus, a descent direction does not exist.

Next, we show that $x$ is not a local minimizer: for any $\varepsilon>0$, select a sufficiently integer large $k$ such that $\alpha =\frac{1}{2k\pi} < \varepsilon$.
Let $\hat{x}$ be any unit vector.
Then we have that $\alpha\hat{x}\in B(\mathbf{0},\varepsilon)$, but also that $f(\alpha \hat{x}) = 0$. 
Hence, $x=\mathbf{0}$ is not a local minimizer.

B.

($\Leftarrow$):
Number the dimension $\{0..n\}$.
Let $e_i$ be the one-hot vector with nonzero coordinate in dimension $i$.
Consider the set of vertices $\{v_i\}_{i\in\{0..n\}}$ where $v_i = \sqrt{x_i} e_i$ and 
\begin{align*}
    v_{0,i} = \begin{cases}
        \sqrt{x_i} &\text{for }i\in \{1..n\}, \\
        \sqrt{1-\sum_{j=1}^n x_j} &\text{for }i=0.
    \end{cases}
\end{align*}
Since $x_i \geq 0, \forall i \in [n]$ and $\sum_{i=1}^n x_i \leq 1$, all entries are real.
Then straightforward calculations verify that $X_{ij} = v_i^\top v_j$, which is an equivalent definition for $X\succeq 0$.

($\Rightarrow$):
If $X\succeq 0$, then all its pricipal submatrices have nonnegative determinants.
In particular, its diagonal entries must be nonnegative, so $x_i\geq 0$. 
Consider the vector $v=e_0-\sum_{i=1}^n e_i$.
Since $X\succeq 0$, we have that
\begin{align*}
    0\leq v^\top X v = 1 - 2 \sum_{i=1}^n x_n + \sum_{i=1}^n x_n.
\end{align*}
Therefore, $\sum_{i=1}^n x_n\leq 1$.


\begin{enumerate}[resume=q]
\item \textbf{Quasiconvexity and Back (10 points)}
\begin{enumerate}[label=\textbf{Part \Alph*.}, leftmargin=*]
\item Prove that $f(x) = \log\bigl(\sum_{i=1}^n a_i e^{x_i}\bigr)$ where $a_i > 0$ is convex.
\item Call a function \emph{kinda-convex} if all its sublevel sets are convex.
\begin{enumerate}[label=\arabic*.]
\item Give a function that is kinda-convex but not convex.
\item Prove that, for any convex $f : \reals \to \reals$ and nondecreasing $g : \reals \to \reals$, $f(g(x))$ is kinda-convex.
\item Can such functions be made convex by stretching the $x$-axis? More precisely, given a kinda-convex $h : \reals \to \reals$, does there exist a differentiable, invertible $g : \reals \to \reals$ such that $h(g(x))$ is convex? Prove your claim.
\end{enumerate}
\end{enumerate}

\textbf{Challenge (Zero Credits):} Try to generalize (A) to $f(X_1, \ldots, X_n) = \log \operatorname{tr}\bigl(\sum_{i=1}^n a_i e^{X_i}\bigr)$. Matrix exponential simply exponentiates the eigenvalues in the spectral decomposition.
\end{enumerate}

\solution

A.

\begin{lemma}[H\"older's inequality]\label{lemma:holder}
    Let $\lambda\in(0,1)$ and $u_i,v_i\geq 0$ for $i\in[n]$. Then
    \[
        \sum_{i=1}^n u_i^{\lambda}v_i^{1-\lambda} \leq \left(\sum_{i=1}^n u_i\right)^{\lambda}\left(\sum_{i=1}^n v_i\right)^{1-\lambda}.
    \]
\end{lemma}
\begin{proof}
    H\"older's inequality states that $\sum_{i=1}^n |s_it_i| \leq \left(\sum_{i=1}^n |s_i|^p\right)^{1/p}\left(\sum_{i=1}^n |t_i|^q\right)^{1/q}$ for all $s,t\in\reals^n$ and all $p,q\in(1,\infty)$ with $\frac{1}{p}+\frac{1}{q}=1$.
    Apply it with $p=\frac{1}{\lambda}$, $q=\frac{1}{1-\lambda}$, $s_i=u_i^{\lambda}$ and $t_i=v_i^{1-\lambda}$, noting that $|s_i|^p=u_i$ and $|t_i|^q=v_i$.
\end{proof}

Since $a_i>0$, we have $\sum_{i=1}^n a_ie^{x_i}>0$ for every $x\in\reals^n$, so $f$ is defined on all of $\reals^n$.
Fix $x,x'\in\reals^n$ and $\lambda\in[0,1]$. If $\lambda\in\{0,1\}$ there is nothing to prove, so let $\lambda\in(0,1)$.
Let $u_i = a_ie^{x_i}$ and $v_i = a_ie^{x_i'}$, which are positive since $a_i>0$. For each $i$,
\[
    a_ie^{\lambda x_i + (1-\lambda)x_i'} = a_i^{\lambda}a_i^{1-\lambda}e^{\lambda x_i}e^{(1-\lambda)x_i'} = u_i^{\lambda}v_i^{1-\lambda}.
\]
Hence, by \cref{lemma:holder} and since $\log$ is increasing,
\begin{align*}
    f(\lambda x + (1-\lambda)x') &= \log\left(\sum_{i=1}^n u_i^{\lambda}v_i^{1-\lambda}\right)\\
    &\leq \log\left(\left(\sum_{i=1}^n u_i\right)^{\lambda}\left(\sum_{i=1}^n v_i\right)^{1-\lambda}\right)\\
    &= \lambda\log\left(\sum_{i=1}^n a_ie^{x_i}\right) + (1-\lambda)\log\left(\sum_{i=1}^n a_ie^{x_i'}\right)\\
    &= \lambda f(x) + (1-\lambda)f(x').
\end{align*}
Thus $f$ is convex.

B.

1. $\sqrt{|x|}$ has convex sublevel sets on $\reals$.

2. If $f$ is convex, the epigraph of $f$ is convex. 
For a fixed $\alpha\in \reals$, the set of points of the epigraph of $f$ below the $y=\alpha$ hyperplane is convex, so the sublevel sets of $f$ are convex, i.e., they are intervals in $\reals$.
Since $g$ is non-decreasing, preimages of intervals are intervals, so $f(g(\cdot))$ is kinda-convex.

3. Consider $h(x) = -e^x$, which is a kinda-convex function, which is also bounded.
Let $f(x) = h(g(x))$. Then $f$ is bounded.
The only bounded convex functions are constants: suppose $f$ is convex and has $x,y\in \reals$ such that $f(x)<f(y)$. 
Let $d=y-x$. Then for any $k>0$, by convexity, $f(x+kd)\geq f(x) + k(f(y)-f(x))$ which is unbounded as $k\to\infty$.
Hence, $f$ is constant. 
Since $h$ is injective, $g$ must be constant, and therefore not invertible.

This is a contradiction, so no such $g$ can exist.

\end{document}
